% Reference LaTeX document for multi-script testing (XeLaTeX)
\documentclass{article}
\usepackage{fontspec}
\usepackage{polyglossia}
\setmainfont{Noto Sans}
\newfontfamily\cjkfont{Noto Sans CJK SC}
\newfontfamily\arabicfont{Noto Sans Arabic}
\newfontfamily\hebrewfont{Noto Sans Hebrew}
\newfontfamily\devanagarifont{Noto Sans Devanagari}
\newfontfamily\tamilfont{Noto Sans Tamil}
\newfontfamily\koreanfont{Noto Sans KR}
\newfontfamily\tibetanfont{Noto Sans Tibetan}
\newfontfamily\ethiopicfont{Noto Sans Ethiopic}

% Add more font families as needed for specific scripts.

\begin{document}
\section*{Multi-Script Reference}

\subsection*{Simple LTR}
Latin: The quick brown fox jumps over the lazy dog.\\
Greek: Καλημέρα κόσμε\\
Cyrillic: Привет мир\\
Georgian: გამარჯობა მსოფლიო\\

\subsection*{Complex Indic}
Devanagari: {\devanagarifont क्योंकि क्षत्रिय}\\
Tamil: {\tamilfont வணக்கம் உலகம்}\\

\subsection*{RTL}
Arabic: {\arabicfont مرحبا بالعالم}\\
Hebrew: {\hebrewfont שלום עולם}\\

\subsection*{CJK}
Chinese: {\cjkfont 你好，世界}\\
Japanese: {\cjkfont こんにちは世界}\\
Korean: {\koreanfont 안녕하세요 세계}\\

\subsection*{Tibetan}
{\tibetanfont བོད་ཀྱི་སྐད་}

\subsection*{African}
Ethiopic: {\ethiopicfont ሰላም ዓለም}

\end{document}
